% !TeX root = main.tex
\chapter{三月二十二日}

うどん屋「な也」。少しばかり高級な味のするうどんをすすり終えた後、私と妻は椅子にもたれてぐったりとしていた。もちろんうどんに食うのに疲れて疲労しているのではない。数時間前に、王子動物園を見て回ったためである。

王子動物園は、パンダを飼育している動物園として名を馳せていた。そのため、最寄りの王子公園駅では壁にパンダのデザインが施されていたり、広報用のウェブサイトでもパンダが中央に陣取って写っている。

パンダの飼育は困難だそうだ。何でも、パンダはあまり交尾をしたがらぬ性質だそうで、なかなか増えてくれない。もちろん、日本に野生のパンダは生息していないため、取って捕まえて増やすこともできない。

そんなわけで、王子動物園のパンダもついに絶えてしまった。

とはいえ、パンダのいない動物園も中々悪くはない。そもそも、私はパンダに特別に愛着があるわけでもないし、ライオン、コアラ、ホッキョクグマなど、人気のある役者はまだ健在だ。

動物博物館というのもあった。動物たちの餌の展示を見てみると、ゴリラは私よりも栄養のある良い飯を食っていた。興味深いものが多数あり大変愉快であったため、動物を見るより良かったと伝えると、妻はしかめっ面をした。

帰り際に妻が、次来るときは子供も一緒に来れたらいいねと言った。本当にそうだと思った。おぼろげな未来を想像しながら、園を後にした。


